\documentclass{article}

\title{A Brief Study on Quantum Entanglement}
\author{John Doe}
\date{June 2025}

\begin{document}
\maketitle

\begin{abstract}
This paper explores the fundamental principles of quantum entanglement and its implications for quantum computing. We present a simplified model for understanding Bell states and their applications in quantum information theory.
\end{abstract}

\section{Introduction}
Quantum entanglement is a phenomenon where two particles become correlated in such a way that the quantum state of one particle cannot be described independently of the other \cite{einstein1935}. This property, famously described by Einstein as ``spooky action at a distance,'' has become a cornerstone of modern quantum information science.

\section{Mathematical Framework}
Consider a two-qubit system. The Bell state $|\Phi^+\rangle$ is defined as:

$$|\Phi^+\rangle = \frac{1}{\sqrt{2}}(|00\rangle + |11\rangle)$$

\begin{theorem}
For any Bell state $|\psi\rangle$, measuring one qubit immediately determines the state of the other, regardless of spatial separation.
\end{theorem}

\begin{proof}
Let $|\psi\rangle = \frac{1}{\sqrt{2}}(|00\rangle + |11\rangle)$. Measuring the first qubit in the computational basis yields $|0\rangle$ with probability $\frac{1}{2}$, collapsing the joint state to $|00\rangle$. Similarly, obtaining $|1\rangle$ collapses the state to $|11\rangle$. In both cases, the second qubit's state is uniquely determined.
\end{proof}

\subsection{Entanglement Entropy}
The von Neumann entropy of the reduced density matrix $\rho_A$ quantifies entanglement:

$$S(\rho_A) = -\text{Tr}(\rho_A \log_2 \rho_A)$$

For a maximally entangled Bell state, $S(\rho_A) = 1$ ebit.

\section{Quantum Circuit Implementation}
We implemented the quantum circuit using standard gates:

\begin{enumerate}
\item Apply Hadamard gate $H$ to qubit 1
\item Apply CNOT gate with qubit 1 as control
\item Measure both qubits
\end{enumerate}

The resulting state after step 2 is:

$$\text{CNOT}(H \otimes I)|00\rangle = \frac{1}{\sqrt{2}}(|00\rangle + |11\rangle)$$

\section{Conclusion}
Quantum entanglement remains one of the most fascinating phenomena in physics. Its applications in quantum computing, quantum cryptography, and quantum teleportation continue to drive research forward.

\begin{thebibliography}{9}
\bibitem{einstein1935} A. Einstein, B. Podolsky, N. Rosen, \textit{Can Quantum-Mechanical Description of Physical Reality Be Considered Complete?}, Physical Review, vol. 47, pp. 777--780, 1935.
\bibitem{bell1964} J. S. Bell, \textit{On the Einstein Podolsky Rosen Paradox}, Physics Physique Fizika, vol. 1, no. 3, pp. 195--200, 1964.
\bibitem{nielsen2010} M. A. Nielsen, I. L. Chuang, \textit{Quantum Computation and Quantum Information}, Cambridge University Press, 2010.
\end{thebibliography}

\end{document}
